% !TeX root = ../colloquium-2.tex

\phantomsection\label{sec_Proofs03}
\begin{proof}[Доказательство критерия непрерывности монотонной функции. Доказательство теоремы об обратной функции.]
    \subsection*{Доказательство критерия непрерывности монотонной функции.}
    \begin{note}[amber]
        Монотонная на отрезке $[a,b]$ функция $f$ непрерывна на этом отрезке тогда и только тогда, когда множеством её значений является отрезок с концами $f(a)$ и $f(b)$.
        Перед доказательством отметим, что $f(a) \le f(b)$, если $f$ не убывает, и $f(a) \ge f(b)$, если $f$ не возрастает. \\
        \begin{note}[purple]
            НЕОБХОДИМОСТЬ. \\
            Если $c \in [a,b]$, то $f(c)$ лежит между $f(a)$ и $f(b)$ в силу монотонности.
            По теореме Коши о непрерывной функции функция $f$ принимает все промежуточные значения между $f(a)$ и $f(b)$.
            Тем самым доказано, что область значений функции $f$ — это отрезок с концами $f(a)$ и $f(b)$.\\
            ДОСТАТОЧНОСТЬ. \\
            Предположим, что функция $f$ монотонна на отрезке $[a,b]$, область её значений — отрезок с концами $f(a)$ и $f(b)$, но $f$ имеет разрыв в точке $x_0 \in [a,b]$.
            Тогда по теореме о разрывах монотонной функции (см. предыдущую лекцию) у $f$ может быть разрыв только первого рода, то есть если $x_0 \in (a,b)$, то в точке разрыва существуют левый и правый пределы, но они не равны между собой:
            \[
            \lim_{x \to x_0 - 0} f(x) = A \ne B = \lim_{x \to x_0 + 0} f(x).
            \]
            Один из интервалов $(A, f(x_0))$ или $(f(x_0), B)$ непуст, и в нём нет значений функции $f$.
            В силу монотонности функции $f$ этот интервал содержится в отрезке с концами $f(a)$ и $f(b)$, поэтому этот отрезок не входит целиком в область значений функции $f$.
            Если же $x_0 = a$, то в этой точке существует правый предел, равный $B \ne f(a)$, а тогда интервал $(f(a), B)$ не содержит ни одного значения функции $f$.
            Случай $x_0 = b$ разбирается аналогично, только теперь в точке $x_0$ существует левый предел $A \ne f(b)$.
            В любом из случаев получаем противоречие с тем, что область значений функции является отрезком.
        \end{note}
    \end{note}

    \subsection*{Доказательство теоремы об обратной функции.}
    \begin{note}[amber]
        Пусть функция $f$ непрерывна и строго монотонна (то есть возрастает или убывает) на отрезке $[a,b]$. Тогда функция $f$ имеет обратную функцию $f^{-1}$, определённую на отрезке с концами $f(a)$ и $f(b)$, 
        причём $f^{-1}$ строго монотонна и непрерывна на отрезке с концами $f(a)$ и $f(b)$, и характер монотонности функций $f$ и $f^{-1}$ одинаковый.\\
        \begin{note}[purple]
           Докозательство. То, что образом функции $f$ является отрезок с концами $f(a)$ и $f(b)$, сразу следует из предыдущей теоремы, то есть $f$ является сюръекцией отрезка $[a,b]$ на отрезок с концами $f(a)$ и $f(b)$.
           Инъекция вытекает из строгой монотонности: двум разным значениям функции соответствуют два разных значения аргумента. Таким образом, $f$ является биекцией, поэтому обратная функция $f^{-1}$ существует по определению.
           Монотонность функции $f^{-1}$ следует из того, что если, например, $f$ возрастает, то
           \[
           f^{-1}(f(x_1)) < f^{-1}(f(x_2)) \iff x_1 < x_2,
           \]
           то есть $f^{-1}$ тоже возрастает, и аналогично для убывания.
           По определению, функция $f^{-1}$ определена на отрезке с концами $f(a)$ и $f(b)$, а областью её значений является отрезок $[a,b]$, поэтому функция $f^{-1}$ непрерывна в силу предыдущей теоремы.
        \end{note}
    \end{note}
\end{proof}
